%%%
%
% $Autor: Wings $
% $Date: 2024-10-31 11:15:45Z $
% $File: file.tex $
% $Version: 1 $
%
% Project: 
%
%%%


\chapter{Schlussfolgerung und Fazit}
\label{chap:Conclusion_Radar}

Das Ziel dieser Projektarbeit, die Konzeption, Konstruktion und informationstechnische Realisierung eines autonomen Ultraschallradarsystems, wurde erfolgreich umgesetzt. Das mechatronische Gesamtsystem demonstriert das Zusammenspiel aus mechanischer Gehäuseauslegung,Hardware und algorithmischer Signalverarbeitung.

Durch die konstruktive Maßnahme der hängenden Servomontage konnte das Problem der Kabel-Kollisionen und Materialbelastung bei rotierenden Systemen weitestgehend gelöst werden. Auf der Softwareseite sorgt das bidirektionale Handshake-Protokoll in Kombination mit dem Runtime-Watchdog für eine Systemstabilität. Die physikalische Schwäche der Ultraschall-Divergenz wurde durch die softwareseitige Hysterese- und Clusterbildung im Processing-Frontend kompensiert, sodass ein Abbild des Nahbereichs angezeigt wird.

Zusammenfassend zeigt das Projekt, dass mit dem Einsatz kostengünstiger Standardkomponenten ein funktionales Messsystem zur digitalen Umfelderfassung realisiert werden kann, welches beispielsweise als Fundament für zukünftige mobile Roboterplattformen dienen kann.